\subsection{The channel results stated}\label{sec:channel-statements}

Four of the rows of Table~\ref{to:tab:corollaries} deserve to be stated as theorems in the main text, because they are the ones a reader is most likely to want to use. In them $\Trt$ is in bits, $B(\Phi)$ is an energy, and the capacity is written in energy units, $\kB T\ln2\,\mathcal T(\Phi)$, wherever it is compared with $B(\Phi)$. Their proofs are in Appendix~\ref{app:channels} (channels) and Appendix~\ref{app:batch} (clocks); the appendix statements carry every hypothesis.

\FloatBarrier
\begin{theorem}[Real-time work law for channels]\label{thm:main-realtime}
Let $\Phi$ be a time-covariant channel on a system with Hamiltonian $H$. In real time, the optimal work per use at vanishing error is $\kB T\ln2\,\Trt(\Phi)$, with
\[
\Trt(\Phi)=\sup_{\rho>0}\Bigl[S(\rho)-S(\Phi(\rho))+m_\rho(\Phi)+\tfrac{\beta}{\ln2}\Tr H\bigl(\Phi(\rho)-\rho\bigr)\Bigr],
\qquad
m_\rho(\Phi)=\inf_{(U,\tau)}I(O{:}E')_\rho,
\]
the infimum over finite exact dilations of the mutual information between the output and the bath output. The thermodynamic capacity is the same expression with $m_\rho=0$. Moreover, for every fixed $\delta>0$, any implementation whose controller starts thermal and that draws at most $n\kB T\ln2\,(\Trt(\Phi)-\delta)$ of deterministic work has error tending to one.
\end{theorem}

\begin{theorem}[Only quantum channels pay]\label{thm:main-dichotomy}
For repeated uses of one memoryless channel, each output released before the next input arrives, as in Theorem~\ref{thm:main-realtime}:
\begin{enumerate}[label=(\alph*),leftmargin=*]
\item If $\Phi$ is classical, meaning that it measures in an eigenbasis of $H$ and prepares states diagonal in that basis, then $\Trt(\Phi)$ equals the capacity at every temperature: waiting saves nothing.
\item For amplitude damping with probability $0<p<1$ on a qubit of gap $\Delta$, $\Trt(\Phi)=\max_{0\le y\le p}[h_2(y)-\beta\Delta y/\ln2]$ exceeds the capacity at every temperature; at $p=1$, a full reset, the two coincide.
\item For the Werner--Holevo channel $W(\rho)=(\Tr(\rho)I-\rho^T)/2$ on a qutrit with zero Hamiltonian, the real-time rate is $\log_23$ bits per use, while in batch the work is $o(n)$; at zero error a batch device too draws at least $n\log_23$ bits of free energy from its battery and controller, a bound attained with an information battery.
\end{enumerate}
\end{theorem}

\begin{theorem}[What a returned clock costs, and how large it must be]\label{thm:main-heisenberg}
Let $\Phi$ have Kraus operators $A_i$, complementary channel $\Psi$, and energy change $D=\Phi^*(H)-H$.
\begin{enumerate}[label=(\alph*),leftmargin=*]
\item \emph{Cost.} With a clock that must be returned, the batch work per use is the single-letter quantity
\[
B(\Phi)=\max_{\Psi(\rho)=\Psi(\sigma)}\Bigl\{\Tr\rho D+\kB T\ln2\bigl[S(\sigma)-S(\Phi(\sigma))\bigr]\Bigr\}\ \ge\ \text{capacity},
\]
with equality whenever $D\in\operatorname{span}\{A_i^*A_j\}$, that is, whenever some environment Hamiltonian can carry the channel's mean energy change; we call such channels balanced. The converse fails: for the Hadamard gate on a qubit of gap $E$, $D\notin\operatorname{span}\{I\}$ and yet $B(\Phi)=\text{capacity}=E/\sqrt2$, because a unitary changes no entropy and both maxima reduce to $\lambda_{\max}(D)$; for $X$-dephasing on the same qubit $B(\Phi)=E/2$ exceeds the capacity strictly.
\item \emph{Size.} When $D$ lies outside the span, the generator $i\sum_iA_i^*\dot A_i$ of the channel's time orbit $\Phi_t(\rho)=U_t\Phi(U_t^*\rho U_t)U_t^*$, $U_t=e^{-itH}$, which equals $D$, satisfies the Hamiltonian-not-in-Kraus-span condition of Zhou and Jiang~\cite{ZhouJiang21}, and the spread of the clock's capital must grow at least linearly in $n$, and at least like $n/\sqrt{\eta_n}$ when the clock is returned within $\eta_n$.
\end{enumerate}
\end{theorem}

\paragraph{Capital and implementation size.} Part~(b) of Theorem~\ref{thm:main-heisenberg} bounds the clock's capital, not the rate: the capital is invested and comes back, and part~(a) gives the work per use whatever its size. The spread is the quantity $s_n$ of Proposition~\ref{kt:prop:capital-converse}. If the controller starts in $\alpha_C=\sum_a\lambda_a|a\rangle\langle a|$ and has Gibbs state $\gamma_C=\sum_b\mu_b|b\rangle\langle b|$, it is the least mean absolute deviation of $\ln(\lambda_a/\mu_b)$ from a constant, with weights $\lambda_a|\langle a|b\rangle|^2$; for a pure controller it is $\beta$ times the mean absolute deviation of its energy from a median. The bound comes from metrology: a code built as that of Zhou and Jiang turns $n$ uses of any reliable implementation into a probe that reads the controller's clock to within $O(1/n)$ (Appendix~\ref{app:batch}). At vanishing error the clock must be broader still, with spread of order at least $n/\eta_n$ when $\varepsilon_n=O(\eta_n)$ (Proposition~\ref{kt:prop:broadclock}). When $D$ lies in the span, a clock of spread $O(\sqrt n)$ suffices for an implementation at fixed error and fixed return error (Proposition~\ref{kt:prop:clocksize}(c)), at work of order $n\log n$; whether such a clock suffices at the optimal rate is open in general (Section~\ref{sec:discussion}), and Proposition~\ref{kt:prop:reach} gives a balanced channel for which it does. For a covariant channel no clock is needed.

\begin{theorem}[Holding is decidable; releasing is not]\label{thm:main-undecidable}
In batch, a channel costs nothing with a returned clock exactly when it is Gibbs-preserving and $D\in\operatorname{span}\{A_i^*A_j\}$, a condition that linear algebra decides for channels and Hamiltonians with algebraic entries when $\beta$ is algebraic, or the Gibbs state $\gamma=e^{-\beta H}/\Tr e^{-\beta H}$ has algebraic entries (for example, every Gibbs weight $e^{-\beta E_k}$ is algebraic). In real time a channel costs nothing exactly when it is a limit of thermal operations, and there is a dimension $d_0$ such that for every $d\ge d_0$, every Hamiltonian diagonal in the dephasing basis and every temperature, no algorithm decides whether a generalized dephasing channel on $\mathbb C^d$ with cyclotomic entries is free in real time; the set of free ones is not recursively enumerable.
\end{theorem}

The batch statement needs a hypothesis on $\beta$ because Gibbs preservation involves the weights $e^{-\beta E_k}$ of the eigenvalues $E_k$ of $H$. If the Gibbs state has algebraic entries, $\Phi(\gamma)=\gamma$ is exact linear algebra over the algebraic numbers. If $\beta$ is algebraic, and nonzero as throughout, the Lindemann--Weierstrass theorem, by which $e^{\alpha_1},\dots,e^{\alpha_m}$ are linearly independent over the algebraic numbers for distinct algebraic $\alpha_k$, turns $\Phi(\gamma)=\gamma$ into the algebraic conditions $\Phi(P_k)=P_k$ for the spectral projections $P_k$ of $H$. The span condition is linear algebra over the algebraic numbers in either case.

At zero Hamiltonian and in units of purity, the rate of Theorem~\ref{thm:main-realtime}, the Werner--Holevo rate of Theorem~\ref{thm:main-dichotomy}(c) and the real-time zero set of Theorem~\ref{thm:main-undecidable} are in our earlier work~\cite{Memory,ImmediateDelivery,Douglas26closed}, and so are a strong converse at zero rate and the undecidability of that zero set when the dimension may grow~\cite{Flagship}. The batch rate of zero in Theorem~\ref{thm:main-dichotomy}(c) is the theorem of Faist, Berta and Brand\~ao, since the capacity of a unital channel at zero Hamiltonian is zero~\cite{FBB21}. Two steps are taken from our earlier preprints and not reproved here. The attainability in Theorem~\ref{thm:main-realtime}, and its formula for $m_\rho$ as an infimum over finite exact dilations, rest on the pointwise smoothing identity~\cite[Thm.~4.1, Cor.~4.3]{Douglas26closed} (Lemma~\ref{sw:lem:smoothing}). The undecidability in Theorem~\ref{thm:main-undecidable} rests on the zero-Hamiltonian closure theorem~\cite[Thm.~3.2(a)]{Memory}, which supplies the direction from zero cost to a limit of finitely factorizable channels (Step~1 of the proof of Theorem~\ref{sw:thm:undecidable}). Theorems~\ref{thm:main-realtime} and~\ref{thm:main-dichotomy} use the signed battery deviation of Table~\ref{tab:exceptions}.

The dichotomy of Theorem~\ref{thm:main-dichotomy} is the first visible consequence of the two layers of Section~\ref{sec:intro}: for a classical channel and an input diagonal in its basis, the bath can be made nearly independent of the output, by a dilation that hides the channel's randomness in a nearly flat register, so $m_\rho=0$ on such inputs. Coherent inputs can leave a positive $m_\rho$ (for complete dephasing in the computational basis at zero Hamiltonian, the input $(I+0.6X)/2$ has $m_\rho=1-h_2(0.8)\approx0.28$ bits), but they cannot raise the optimized real-time cost, so the real-time rate is the capacity (Appendix~\ref{app:channels}). For a quantum channel such as partial reset the output is necessarily correlated with what is left behind. The dichotomy concerns memoryless channels: a classical process with memory can pay for release, as the task of Figure~\ref{fig:three} does (Section~\ref{sec:classes}). Theorem~\ref{thm:main-heisenberg} says that the clock-size question and the metrology question have the same answer for the same reason: inputs that the environment cannot tell apart but that differ in energy are exactly the inputs that make a code sensitive to the channel's phase, as in the code of Zhou and Jiang~\cite[Sec.~V]{ZhouJiang21}. Whether the clock also costs work above the capacity is a separate question, decided by $B(\Phi)$ and not by the span alone.
